	 
	 \newpage
 % =========================================================
	 \subsection{LR-Zerlegung} \label{chap:lr_zerlegung}
 % =========================================================
	
	\defi{ \n	
		Der LR-Zerlegungs Algorithmus soll eine bessere Variante sein, ein System der Form $Ax = b$ auszurechnen. Dazu trennt es die Matrix in zwei andere Matrizen $L$ und $R$ auf, die weitere Berechnungen einfacher gestalten lassen sollen.
		\n
		Was ist $L$, was ist $R$?
		\n
		Man schaue sich die Gauß Zeilenoperationen an einem allgemeinen Beispiel genauer an:
		\begin{align*}
			\begin{bmatrix}
				a & d & g\\
				b & e & h\\
				c & f & i
			\end{bmatrix} \Longrightarrow
			\begin{bmatrix}
				a & d & g\\
				0 & e - \frac{d \cdot b}{a} & h - \frac{g \cdot  b}{a}\\
				c & f & i
			\end{bmatrix} \Longrightarrow \dots
		\end{align*}
		Nun ist die Idee, die Umformung jedes Schrittes nicht als Zeilenoperation umzusetzen, sondern eine Matrix zu finden, die multipliziert mit $A$ dasselbe Ergebnis liefert.
%	}
%	
		\n	
		Obigen Umformungsschritt kann man auch folgendermaßen schreiben:
		\begin{align*}
			\begin{bmatrix}
				1 & 0 & 0\\
				-\frac{b}{a} & 1 & 0\\
				0 & 0 & 1
			\end{bmatrix} \cdot
			\begin{bmatrix}
				a & d & g\\
				b & e & h\\
				c & f & i
			\end{bmatrix} = 
			\begin{bmatrix}
			a & d & g\\
			0 & e - \frac{d \cdot b}{a} & h - \frac{g \cdot  b}{a}\\
			c & f & i
			\end{bmatrix}
		\end{align*}
		Wenn man mehrere Beispiele aufschreiben würde, kann man sehen, dass alle konstruierten Matrizen ähnlich der obigen ist. Wichtige Werte wie bei dem Beispiel das $-\frac{b}{a}$ stehen immer im links unteren Bereich der Matrix und die Diagonaleinträge sind immer Einsen.
%	}
%	
		\n 
		Die Idee der LR-Zerlegung ist es jetzt, alle diese konstruierten Matrizen aufeinander zu multiplizieren; wir nennen sie die inverse $L$ Matrix, also $L\inv$. Nun beschreibt $R$ die Matrix, die man bekommt, wenn man eine reguläre Vorwärtssubstitution ausführen würde, also eine rechts obere Dreiecksmatrix (wie wenn man Gauß Elimination anwenden würde).
		Nochmal in kurz: 
		\begin{align} 
			L\inv \cdot A = R \sepArrow A = L \cdot R
		\end{align}
		Zusammengefasst ist $L$ eine links untere Dreiecksmatrix mit Einsen auf der Diagonalen und $R$ eine rechts obere Dreiecksmatrix.
		\n
		Nochmal zur Klarstellung: $A = L \cdot R$. Daher auch der Name LR-Zerlegung. Im Englischen spricht man von der LU-Zerlegung, da dort nicht links und rechts namensgebend war, sondern die untere (lower) $L$ und die obere (upper) $U$ Dreiecksmatrix. Es ist aber genau dasselbe!
	}
	
	\sbreak
	\methode{ Zerlegung der Matrix in $L$ und $R$: \n
		Herleitungen sind toll, bringen uns aber auch nicht viel weiter, wenn wir nicht wissen, wie man $L$ und $R$ überhaupt ausrechnet.
		\n
		Es gibt mehrere Methoden, die Matrix zu zerlegen und $L$ und $R$ zu finden. Ich stelle euch jetzt meinen persönlichen Favoriten vor (das bedeutet nicht, dass es für euch unbedingt die beste Methode ist!).
		\n
		Schreiben wir auf, was wir über die Matrizen wissen:
		\begin{align}
			\begin{bmatrix}
			& & \\
			& A& \\
			& & 
			\end{bmatrix} = 
			\begin{bmatrix}
			1 & 0 & 0\\
			& 1 & 0\\
			&   & 1
			\end{bmatrix} \cdot
			\begin{bmatrix}
			&&&\\
			0&&\\
			0&0&
			\end{bmatrix}
		\end{align}
		Diese Gleichung lässt sich lösen, da $A$ bekannt ist und wir mittels Matrixmultiplikation auf die Werte in $L$ und $R$ schließen können.
	}
	
	\sbreak
	\bsp{ \n Achtung, sehr lang!
		\begin{align*}
			\begin{bmatrix}
				1& 2& 2\\
				2& 1& -2\\
				3& 0& 2 
			\end{bmatrix}
			 = 
			\begin{bmatrix}
				1 & 0 & 0\\
				& 1 & 0\\
				&   & 1
			\end{bmatrix} \cdot
			\begin{bmatrix}
				  &	  &\\
				0 &   &\\
				0 & 0 &
			\end{bmatrix}
		\end{align*}
		Schauen wir uns den ersten Wert an:
		\begin{align*}
			\begin{bmatrix}
				\green{1}& 2& 2\\
				2& 1& -2\\
				3& 0& 2 
			\end{bmatrix} &= 
			\begin{bmatrix}
				\green{1} & \green{0} & \green{0}\\
				& 1 & 0\\
				&   & 1
			\end{bmatrix} 
			\cdot
			\begin{bmatrix}
				\green{x}&&&\\
				\green{0}&&\\
				\green{0}&0&
			\end{bmatrix} \\
	 		1\cdot x + 0\cdot0 &+ 0\cdot0 = 1 \quad \Longrightarrow x = 1 \\
			\begin{bmatrix}
				1& 2& 2\\
				2& 1& -2\\
				3& 0& 2 
			\end{bmatrix} &= 
			\begin{bmatrix}
				1 & 0 & 0\\
				& 1 & 0\\
				&   & 1
			\end{bmatrix} \cdot
			\begin{bmatrix}
				\red{1}&&&\\
				0&&\\
				0&0&
			\end{bmatrix}
		\end{align*}
		%$$ \sbreak
		Alle weiteren Schritte analog (etwas schneller)
		\begin{align*}
			\begin{bmatrix}
				1& 2& 2\\
				2& 1& -2\\
				3& 0& 2 
			\end{bmatrix} &= 
			\begin{bmatrix}
				1 & 0 & 0\\
				\red{2}& 1 & 0\\
				&   & 1
			\end{bmatrix} \cdot
			\begin{bmatrix}
				1 & &&\\
				0 & &\\
				0 & 0&
			\end{bmatrix} \\
			\text{weil } x\cdot 1 + 1\cdot0 &+ 0\cdot0 = 2 \quad \Longrightarrow x = 2
		\end{align*}
		\n 
		\begin{align*}
			\begin{bmatrix}
				1& 2& 2\\
				2& 1& -2\\
				3& 0& 2 
			\end{bmatrix} &= 
			\begin{bmatrix}
				1 & 0 & 0\\
				2& 1 & 0\\
				&  & 1
			\end{bmatrix} \cdot
			\begin{bmatrix}
				1& \red{2}&& \\
				0 & &\\
				0 & 0 &
			\end{bmatrix} \\
			\text{weil } 1\cdot x + 0 &+ 0\cdot 0 = 2 \quad \Longrightarrow x = 2 \\
			\begin{bmatrix}
				1& 2& 2\\
				2& 1& -2\\
				3& 0& 2 
			\end{bmatrix} &= 
			\begin{bmatrix}
				1 & 0 & 0\\
				2& 1 & 0\\
				&  & 1
			\end{bmatrix} \cdot
			\begin{bmatrix}
				1& 2& \\
				0 & \red{-3}&\\
				0 & 0 &
			\end{bmatrix} \\
			\begin{bmatrix}
				1& 2& 2\\
				2& 1& -2\\
				3& 0& 2 
			\end{bmatrix} &= 
			\begin{bmatrix}
				1 & 0 & 0\\
				2& 1 & 0\\
				\red{3}&  & 1
			\end{bmatrix} \cdot
			\begin{bmatrix}
				1& 2& \\
				0 & -3&\\
				0 & 0 &
			\end{bmatrix} \\
			\begin{bmatrix}
				1& 2& 2\\
				2& 1& -2\\
				3& 0& 2 
			\end{bmatrix} &= 
			\begin{bmatrix}
				1 & 0 & 0\\
				2& 1 & 0\\
				3& \red{2} & 1
			\end{bmatrix} \cdot
			\begin{bmatrix}
				1& 2& \\
				0 & -3&\\
				0 & 0 &
			\end{bmatrix} \\
			\begin{bmatrix}
				1& 2& 2\\
				2& 1& -2\\
				3& 0& 2 
			\end{bmatrix} &= 
			\begin{bmatrix}
				1 & 0 & 0\\
				2& 1 & 0\\
				3& 2 & 1
			\end{bmatrix} \cdot
			\begin{bmatrix}
				1& 2& \red{2}\\
				0 & -3&\\
				0 & 0 &
			\end{bmatrix} \\
			\begin{bmatrix}
				1& 2& 2\\
				2& 1& -2\\
				3& 0& 2 
			\end{bmatrix} &= 
			\begin{bmatrix}
				1 & 0 & 0\\
				2& 1 & 0\\
				3& 2 & 1
			\end{bmatrix} \cdot
			\begin{bmatrix}
				1& 2& 2\\
				0 & -3& \red{-6}\\
				0 & 0 & \red{8}
			\end{bmatrix}
		\end{align*}
		\n
		Somit ergibt sich:
		\begin{align*}
			L =
			\begin{bmatrix}
				1 & 0 & 0\\
				2& 1 & 0\\
				3& 2 & 1
			\end{bmatrix} \quad \quad
			R = 
			\begin{bmatrix}
				1& 2& 2\\
				0 & -3& -6\\
				0 & 0 & 8
			\end{bmatrix}
		\end{align*}
	}
	
	\sbreak
	\defi{	\fat{Vorwärts- und Rückwärtssubst.}
		\n 
		Wenn wir jetzt die Vorwärts- und Rückwärtssubst. noch modifizieren, erhalten wir ein Verfahren, welches ein beliebiges LGS lösen kann.
		\n 
		Also vom Prinzip: 
		\begin{align*} 
			Ax &= b \\
			A &= L \cdot R \\
			\sepArrow LRx &= b 
		\end{align*} 
		Substituieren wir $Rx$ mit $y$, erhalten wir nachfolgende Methode.
	}

	\sbreak
	\methode{
		\n
		LR-Zerlegungs-Verfahren zum Lösen eines LGS der Form $Ax = b$:
		\begin{enumerate}
			\item Zerlegung der Matrix: 
				\begin{align}
					A = L \cdot R
				\end{align}
			\item Vorwärtssubstitution: 
				\begin{align}
					Ly = b
				\end{align}
			\item Rückwärtssubstitution: 
				\begin{align}
					Rx = y
				\end{align}
		\end{enumerate}
	}

	\sbreak
	\bsp{
		\n 
		Lösen wir vorheriges Beispiel zuende. Wir nehmen als $b$ dasselbe wie beim Gauß-Beispiel:
		\begin{align*}
			L =
			\begin{bmatrix}
				1 & 0 & 0\\
				2& 1 & 0\\
				3& 2 & 1
			\end{bmatrix} \quad \quad
			R = 
			\begin{bmatrix}
				1& 2& 2\\
				0 & -3& -6\\
				0 & 0 & 8
			\end{bmatrix} \quad \quad 
			b = 
			\nvector{
				3 \\ 2 \\ 6
			}
		\end{align*}
		Vorwärtssubstitution, also $Ly = b$:
		\begin{align*}
			\begin{bmatrix}
				1 & 0 & 0\\
				2& 1 & 0\\
				3& 2 & 1
			\end{bmatrix} \cdot
			\nvector{
				y_0 \\ y_1 \\ y_2
			}
			= 
			\nvector{
				3 \\ 2 \\ 6
			}
		\end{align*}
		Lösen, indem wir nacheinander von oben nach unten einsetzen:
		\begin{align*}
			1 y_0 &= 3 \\
			& \\
			2 y_0 + 1 y_1 &= 2 \\
			y_1 &= -4 \\
			& \\
			3 y_0 + 2 y_1 + y_2 &= 6 \\
			9 + -8 + y_2 &= 6 \\
			y_2 &= 5
		\end{align*}
		\n
		Rückwärtssubstitution, also $Rx = y$:
		\begin{align*}
			\begin{bmatrix}
				1& 2& 2\\
				0 & -3& -6\\
				0 & 0 & 8
			\end{bmatrix}
			\cdot 
			\nvector{
				x_0 \\ x_1 \\ x_2
			}
			= 
			\nvector{
				3 \\ -4 \\ 5
			}
		\end{align*}
		Jetzt haben wir wieder genau dasselbe System wie nach Gauß, die Lösung ist logischerweise dieselbe (siehe \refChapter{chap:gausselim}).
	}

	\sbreak
	\defi{
		\n
		Vorwärts und Rückwärtssubstitution aus dem Algorithmus sind dann auszurechnende LGSe mit der schönen Eigenschaft, dass man sie verhältnismäßig einfach ausrechnen kann (wegen den Dreiecksmatrizen).
		\n 
		Man kann jetzt sehr schnell sehen, warum das so effizient sein kann.
		Denn egal was $b$ in der Gleichung ist, die Zerlegung für $A$ bleibt dieselbe und muss nach dem ersten Mal nicht mehr ausgerechnet werden.
		\n 
		So liegt der Aufwand bei der Berechnung mit Zerlegung in $\bigO(n^3)$, ohne Berechnung der Zerlegung (wenn wir sie schon haben) nur in $\bigO(n^2)$ und damit schneller als Gauß-Elimination.
	}
	
	
	\sbreak
	\wichtig{
		\n
		Die LR-Zerlegung ist eine schöne Anwendung, wenn wir mit derselben Matrix $A$ viele LGS der Form $Ax = b$ lösen müssen!
	}

